\documentclass[a4paper,landscape]{article}

% Черновик графа формул. Источник — docs/card-titles.tex.

\usepackage{cmap}
\usepackage[T2A]{fontenc}
\usepackage[utf8]{inputenc}
\usepackage[russian]{babel}
\usepackage{amsmath,amssymb}
\usepackage{tempora}
\usepackage[a4paper,landscape,left=12mm,right=12mm,top=12mm,bottom=12mm]{geometry}
\usepackage{array,booktabs,longtable}
\usepackage{ragged2e}
\usepackage[table]{xcolor}

\definecolor{rule}{HTML}{C9CDC7}
\arrayrulecolor{rule}
\renewcommand{\arraystretch}{1.25}
\setlength{\parindent}{0pt}
\setlength{\tabcolsep}{0pt}

\begin{document}

{\LARGE\bfseries Молекулярная физика}\par
\vspace{1mm}
{\large Черновик графа формул}\par
\vspace{2mm}
{\footnotesize
\textbf{Тип}: определение — язык и определения величин; основа — принимаем без вывода;
следствие — выводится из карточек с номерами из столбца «Из чего».
У определения в шагах может быть равносильная запись, у основы — подсказка.
Шаги — черновик, правьте формулы и текст. Номера — из первого столбца.
Всего карточек: 41.}\par
\vspace{4mm}

\begin{longtable}{@{}>{\RaggedRight\arraybackslash}p{7mm}@{\hspace{3mm}}>{\RaggedRight\arraybackslash}p{58mm}@{\hspace{3mm}}>{\RaggedRight\hyphenpenalty=10000\arraybackslash}p{58mm}@{\hspace{3mm}}>{\RaggedRight\arraybackslash}p{24mm}@{\hspace{3mm}}>{\RaggedRight\arraybackslash}p{35mm}@{\hspace{3mm}}>{\RaggedRight\arraybackslash}p{76mm}@{}}
\toprule
\textbf{№} & \textbf{Формула} & \textbf{Название} & \textbf{Тип} & \textbf{Из чего} & \textbf{Шаги} \\
\midrule
\endfirsthead
\toprule
\textbf{№} & \textbf{Формула} & \textbf{Название} & \textbf{Тип} & \textbf{Из чего} & \textbf{Шаги} \\
\midrule
\endhead
\bottomrule
\endfoot
1 & $\displaystyle \nu=\dfrac{N}{N_A}$ & Количество вещества (через число частиц) & определение & — & — \\
\midrule
2 & $\displaystyle M=m_0 N_A$ & Молярная масса (через массу молекулы) & определение & — & — \\
\midrule
3 & $\displaystyle m=m_0 N$ & Масса газа (через массу молекулы) & определение & — & — \\
\midrule
4 & $\displaystyle n=\dfrac{N}{V}$ & Концентрация молекул газа & определение & — & — \\
\midrule
5 & $\displaystyle \nu=\dfrac{m}{M}$ & Количество вещества (через массу тела) & следствие & 1, 2, 3 & $M=m_0N_A$ (2), $m=m_0N$ (3)\newline $\Rightarrow\ \dfrac{m}{M}=\dfrac{N}{N_A}=\nu$ (1) \\
\midrule
6 & $\displaystyle m_0=\dfrac{m}{N}=\dfrac{M}{N_A}$ & Масса одной молекулы & следствие & 2, 3 & $m=m_0N$ (3) $\Rightarrow\ m_0=\dfrac{m}{N}$\newline $M=m_0N_A$ (2) $\Rightarrow\ m_0=\dfrac{M}{N_A}$ \\
\midrule
7 & $\displaystyle N=\nu N_A=\dfrac{m}{M}N_A$ & Число молекул & следствие & 1, 5 & $\nu=\dfrac{N}{N_A}=\dfrac{m}{M}$ (1, 5)\newline $\Rightarrow\ N=\nu N_A=\dfrac{m}{M}N_A$ \\
\midrule
8 & $\displaystyle n=\dfrac{N}{V}=\dfrac{\nu N_A}{V}$ & Концентрация молекул (через количество вещества) & следствие & 4, 1 & $N=\nu N_A$ (1), $n=\dfrac{N}{V}$ (4)\newline $\Rightarrow\ n=\dfrac{\nu N_A}{V}$ \\
\midrule
9 & $\displaystyle \rho=\dfrac{m}{V}$ & Плотность вещества & определение & — & — \\
\midrule
10 & $\displaystyle \rho=m_0 n$ & Плотность газа (через концентрацию молекул) & следствие & 3, 4 & $m=m_0N$ (3), $n=\dfrac{N}{V}$ (4) $\Rightarrow\ \rho=\dfrac{m}{V}=m_0n$ \\
\midrule
11 & $\displaystyle p=\dfrac{1}{3}n m_0\langle v^2\rangle$ & Основное уравнение МКТ & основа & — & — \\
\midrule
12 & $\displaystyle \langle E_{\text{к}}\rangle=\dfrac{m_0\langle v^2\rangle}{2}$ & Средняя кинетическая энергия поступательного движения молекулы & определение & — & — \\
\midrule
13 & $\displaystyle p=\dfrac{2}{3}n \langle E_{\text{к}}\rangle$ & Давление газа (через среднюю кинетическую энергию молекулы) & следствие & 11, 12 & $p=\dfrac{1}{3}nm_0\langle v^2\rangle$ (11), $m_0\langle v^2\rangle=2\langle E_{\text{к}}\rangle$ (12)\newline $\Rightarrow\ p=\dfrac{2}{3}n\langle E_{\text{к}}\rangle$ \\
\midrule
14 & $\displaystyle T=t+273$ & Перевод температуры из шкалы Цельсия в Кельвина & определение & — & — \\
\midrule
15 & $\displaystyle p=nkT$ & Давление газа (через концентрацию и температуру) & основа & — & — \\
\midrule
16 & $\displaystyle \langle E_{\text{к}}\rangle=\dfrac{3}{2}kT$ & Средняя кинетическая энергия молекулы (через температуру) & основа & — & — \\
\midrule
17 & $\displaystyle R =k N_A = 8,31 $ Дж/(моль$\cdot$K) & Универсальная газовая постоянная & определение & — & — \\
\midrule
18 & $\displaystyle v_{\text{кв}}=\sqrt{\langle v^2\rangle}$ & Средняя квадратичная скорость & определение & — & — \\
\midrule
19 & $\displaystyle v_{\text{кв}}=\sqrt{\dfrac{3kT}{m_0}}$ & Средняя квадратичная скорость (через температуру) & следствие & 12, 16, 18 & $\langle E_{\text{к}}\rangle=\dfrac{m_0\langle v^2\rangle}{2}$ (12), $\langle E_{\text{к}}\rangle=\dfrac{3}{2}kT$ (16)\newline $\Rightarrow\ \langle v^2\rangle=\dfrac{3kT}{m_0}$\newline $v_{\text{кв}}=\sqrt{\langle v^2\rangle}=\sqrt{\dfrac{3kT}{m_0}}$ (18) \\
\midrule
20 & $\displaystyle v_{\text{кв}}=\sqrt{\dfrac{3RT}{M}}$ & Средняя квадратичная скорость (через молярную массу) & следствие & 19, 17, 6 & $k=\dfrac{R}{N_A}$ (17), $m_0=\dfrac{M}{N_A}$ (6) $\Rightarrow\ \dfrac{3kT}{m_0}=\dfrac{3RT}{M}$\newline $v_{\text{кв}}=\sqrt{\dfrac{3RT}{M}}$ \\
\midrule
21 & $\displaystyle \dfrac{pV}{T}=\text{const}$ & Уравнение Клапейрона & основа & — & — \\
\midrule
22 & $\displaystyle pV=\dfrac{m}{M}RT=\nu RT$ & Уравнение Клапейрона—Менделеева & основа & — & — \\
\midrule
23 & $\displaystyle \dfrac{p_1V_1}{T_1}=\dfrac{p_2V_2}{T_2}$ & Уравнение Клапейрона для двух состояний & следствие & 21 & $\dfrac{pV}{T}=\text{const}$ (21) $\Rightarrow\ \dfrac{p_1V_1}{T_1}=\dfrac{p_2V_2}{T_2}$ \\
\midrule
24 & $\displaystyle p_{\text{смеси}}=p_1+p_2+\dots$ & Закон Дальтона & основа & — & — \\
\midrule
25 & $\displaystyle p=\dfrac{\rho RT}{M}$ & Давление (через плотность) & следствие & 22, 9 & $pV=\dfrac{m}{M}RT$ (22) $\Rightarrow\ p=\dfrac{m}{V}\dfrac{RT}{M}$\newline $\rho=\dfrac{m}{V}$ (9) $\Rightarrow\ p=\dfrac{\rho RT}{M}$ \\
\midrule
26 & $\displaystyle pV=\text{const}$ & Закон Бойля—Мариотта (изотермический) & следствие & 21 & $\dfrac{pV}{T}=\text{const}$ (21), $T=\text{const}$ $\Rightarrow\ pV=\text{const}$ \\
\midrule
27 & $\displaystyle \dfrac{V}{T}=\text{const}$ & Закон Гей-Люссака (изобарный) & следствие & 21 & $\dfrac{pV}{T}=\text{const}$ (21), $p=\text{const}$ $\Rightarrow\ \dfrac{V}{T}=\text{const}$ \\
\midrule
28 & $\displaystyle \dfrac{p}{T}=\text{const}$ & Закон Шарля (изохорный) & следствие & 21 & $\dfrac{pV}{T}=\text{const}$ (21), $V=\text{const}$ $\Rightarrow\ \dfrac{p}{T}=\text{const}$ \\
\midrule
29 & $\displaystyle U=\dfrac{3}{2}\dfrac{m}{M}RT=\dfrac{3}{2}pV$ & Внутренняя энергия одноатомного идеального газа & основа & — & Подсказка: $U=N\langle E_{\text{к}}\rangle$ \\
\midrule
30 & $\displaystyle \Delta U=\dfrac{3}{2}\dfrac{m}{M} R\Delta T$ & Изменение внутренней энергии идеального газа & следствие & 29 & $U=\dfrac{3}{2}\dfrac{m}{M}RT$ (29) $\Rightarrow\ \Delta U=\dfrac{3}{2}\dfrac{m}{M}R\Delta T$ \\
\midrule
31 & $\displaystyle Q=\Delta U+A=\Delta U+p\Delta V$ & Первый закон термодинамики & основа & — & — \\
\midrule
32 & $\displaystyle A=p\Delta V=p(V_2-V_1)$ & Работа газа при постоянном давлении & определение & — & — \\
\midrule
33 & $\displaystyle Q=cm\Delta T=cm(T_2-T_1)$ & Количество теплоты (нагревание и охлаждение) & определение & — & — \\
\midrule
34 & $\displaystyle Q=\lambda m$ & Количество теплоты (плавление и кристаллизация) & определение & — & Плавление: $Q=\lambda m$; кристаллизация: $Q=-\lambda m$ \\
\midrule
35 & $\displaystyle Q=Lm$ & Количество теплоты (парообразование и конденсация) & определение & — & Парообразование: $Q=Lm$;\newline конденсация: $Q=-Lm$ \\
\midrule
36 & $\displaystyle Q=qm$ & Количество теплоты (сгорание топлива) & определение & — & — \\
\midrule
37 & $\displaystyle \eta=\dfrac{A_{\text{ц}}}{Q_1}$ & КПД теплового двигателя & определение & — & — \\
\midrule
38 & $\displaystyle A_{\text{ц}} = Q_1 - |Q_2|$ & Работа, совершенная двигателем & следствие & 31 & $Q=\Delta U+A$ (31), цикл $\Delta U=0$ $\Rightarrow\ Q=A$\newline $Q=Q_1-|Q_2|$ $\Rightarrow\ A_{\text{ц}}=Q_1-|Q_2|$ \\
\midrule
39 & $\displaystyle \eta=\dfrac{Q_1-|Q_2|}{Q_1}$ & КПД теплового двигателя (через количество теплоты) & следствие & 37, 38 & $\eta=\dfrac{A_{\text{ц}}}{Q_1}$ (37), $A_{\text{ц}}=Q_1-|Q_2|$ (38)\newline $\Rightarrow\ \eta=\dfrac{Q_1-|Q_2|}{Q_1}$ \\
\midrule
40 & $\displaystyle \varphi=\dfrac{p_{\text{п}}}{p_{\text{н}}}\cdot100\%$ & Относительная влажность через давление & определение & — & — \\
\midrule
41 & $\displaystyle \varphi=\dfrac{\rho_{\text{п}}}{\rho_{\text{н}}}\cdot100\%$ & Относительная влажность через плотность & следствие & 40, 25 & $\varphi=\dfrac{p_{\text{п}}}{p_{\text{н}}}$ (40), $p=\dfrac{\rho RT}{M}$ (25) при одной $T$\newline $\Rightarrow\ \varphi=\dfrac{\rho_{\text{п}}}{\rho_{\text{н}}}$ \\
\end{longtable}

\end{document}
